\section{Financial Markets}

In a previous chapter we provided the definition of an asset as any resource owned or controlled by an economic entity
that can be used to produce positive economic value (\ref{def:asset}). We grouped assets into three categories:
tangible assets, intangible assets, and financial assets, with the latter being defined as non-physical assets whose
value is derived from a contract (\ref{def:financial_asset}). This chapter will be focused specificaly on financial
assets. \v

First, let us repeat the definition of a financial asset.

\bd[Financial Asset]
A \textbf{financial asset} is a non-physical asset whose value is derived from a contract.
\ed

The contract that specifies the value of the financial asset is called a ``financial instrument''.

\bd[Financial Instrument]
A \textbf{financial instrument} is a contract that gives rise to a financial asset.
\ed

Sometimes the terms ``financial asset'' and ``financial instrument'' are used interchangeably, but it is important to
understand the difference between them. A financial asset is the result of a financial instrument, and a financial
instrument is a contract that gives rise to a financial asset. The financial instrument is the cause and the financial
asset is the effect. \v

In a previous chapter we also provided the definition of a market as a place where buyers and sellers come together to
trade goods and services (\ref{def:market}). Having defined financial instruments, we can now define a very special kind
of market called ``financial market''.

\bd[Financial Market]
A \textbf{financial market} is a market where people trade financial instruments at low transaction costs and at prices
that reflect supply and demand.
\ed

Financial markets can be divided into two categories: ``money markets'' and ``capital markets''.

\bd[Money Market]
A \textbf{money market} is a financial market in which short-term (a year or less) financial instruments are traded.
\ed

The most important financial instruments that are traded in money markets are treasury bills, commercial papers,
certificates of deposit, and others.

\bd[Capital Market]
A \textbf{capital market} is a financial market in which long-term (more than a year) financial instruments are traded.
\ed

The most important financial instruments that are traded in capital markets are securities, derivatives, currencies and
insurances. \v

For the rest of this chapter we will focus specifically on capital markets, and the financial instruments that are
traded in them. But before we do that, we will provide a brief overview of the main economic functions of capital
markets which are:
\bit
\item \textbf{Wealth Allocation} To facilitate the raising of capital by channeling wealth and allocating savings into
investments.
\item \textbf{Liquidity}: To allow the flexibility in timing consumption by providing liquidity of buying and selling
financial instruments.
\item \textbf{Risk Management}: To provide a mechanism for risk management including risk transfer, risk reduction
through diversification, and optimal risk sharing.
\item \textbf{Pricing}: To provide a mechanism for price discovery, i.e.\ the process of determining the price of an
asset in the market.
\eit

Hence, any time one engages in a capital market transaction they're altering the risk profile of their asset portfolios.
Subsequently, another way of thinking about capital markets is not necessarily the movement of capital, but the
alteration of those risks.

\section{Securities Markets}

\bd[Security]
A \textbf{security} is a financial instrument that represents an ownership in a corporation or a creditor relationship
with a governmental body or a corporation.
\ed

\bd[Securities Market]
A \textbf{securities market} is a market where people trade securities.
\ed

The main function of security markets is wealth allocation, i.e.\ to facilitate the transfer of funds from those who
have excess funds to those who need funds, by providing a mechanism for the buying and selling of securities.

\subsection{Market Participants}

\bd[Investor]
An \textbf{investor} is a person or entity that commits capital with the expectation of receiving financial returns
generated by securities.
\ed

In traditional finance investors are considered to be just households, however in modern finance an investor can be any
institution who has capital to invest including pension funds investing group retirement plans, insurance companies
investing premiums, investment companies like mutual funds and unit trusts, sovereign wealth funds, endowments like
universities and non-profit oranizations, asset managers, and hedge funds. Its worths mentioning though that all these
institutional investors are ultimately investing the savings of households, hence in a way they act in behalf of the
households.

\bd[Issuer]
An \textbf{issuer} is a legal entity that develops, registers and sells securities for the purpose of financing its
operations.
\ed

In traditional finance issuers are considered to be just corporations, however in modern finance an issuer can be any
entity who needs to raise capital, including governments and governments agencies, municipalities, and supranational
entities. \v

Although issuers are the one who issue securities, most of the time they do not sell them directly to investors.
Instead,they hire financial intermediaries between them and the investors.

\bd[Financial Intermediary]
A \textbf{financial intermediary} is an institution that acts as the middleman between investors and issuers.
\ed

The main function of financial intermediaries is to facilitate the transfer of funds from those who have excess funds
to those who need funds, by providing a mechanism for the buying and selling of securities. The most important
financial intermediaries are comercial and investment banks.

\fig{cm}{0.5}

\bd[Commercial Bank]
A \textbf{commercial bank} is a financial institution that provides deposit taking and commercial lending services.
\ed

\fig{cm1}{0.5}

\bd[Investment Bank]
An \textbf{investment bank} is a financial institution that provides securities trading, securities underwriting and
corporate advisory services.
\ed

\fig{cm2}{0.5}

Althoug there is a clear distinction between commercial and investment banks, in practice the distinction is not always
clear. In the past, commercial banks were not allowed to engage in investment banking activities, and investment banks
were not allowed to engage in commercial banking activities. However, in recent years, the distinction between the two
types of banks has become less clear, finding the two types of banks competing for similar services, and many banks
providing both types of services. \v

Securities markets can be divided into two categories based on the nature of the transactions: ``primary markets'' and
``secondary markets''.

\subsection{Primary Market}

\bd[Primary Market]
A \textbf{primary market} is a market where new securities are issued and sold for the first time.
\ed

Primary markets are where new securities are issued and sold for the first time through initial public offerings (IPOs)
or private placements.

\bd[Initial Public Offering (IPO)]
An \textbf{initial public offering} (\textbf{IPO}) is the first time that a company's securities are offered to the
public.
\ed

\bd[Private Placement]
A \textbf{private placement} is a sale of securities to pre-selected investors and institutions rather than on the
open market.
\ed

\subsection{Secondary Market}

\bd[Secondary Market]
A \textbf{secondary market} is a market where existing securities are bought and sold.
\ed

Secondary market is the trading of already outstanding securities. For the most part, they are non issuer
transactions. That is, the trade takes place between two different investors that are exchanging money for securities,
but it doesn't change the number of securities outstanding, nor does it provide the issuer with any additional
capital for investment. In simple words, there's no new money being raised and no wealth being channeled into
investment in the secondary market, however secondary markets are much larger than the primary markets, with the vast
majority of securities transactions (more than 95\%) taking place in the secondary market. \v

Although there are no issuers in secondary markets, there are financial intermediaries. The two primary roles of
financial intermediaries in the secondary market are to act either as ``market makers'' (also called ``dealers'') or
as ``brokers''.

\bd[Market Maker / Dealer]
A \textbf{market maker} or \textbf{dealer} is a financial intermediary that acts as principal by standing ready to
trade securities at all times at the publicly quoted price.
\ed

In simple words, a market maker is a financial intermediary in the business of trading securities for their own
account, and not on behalf of clients. Market makers take the opposite side of a trade from the customer, i.e.\ they
buy when a customer wants to sell and sell when a customer wants to buy. As it makes sense, this makes market makers
to have some capital in risk, due to the security inventory they have, and they need to be compensated for that. In
order to do so, market makers quote their own prices on securities both for buying (called ``bid price'') and for
selling (called ``ask price'' or ``offer price'').

\bd[Bid Price]
The \textbf{bid price} is the price at which a market maker is willing to buy a security.
\ed

\bd[Ask Price / Offer Price]
The \textbf{ask price}, or \textbf{offer price}, is the price at which a market maker is willing to sell a security.
\ed

It is important to note that both bid and ask prices do not reflect the true value of a security, but rather the market
maker's opinion of the security's value. The difference between the bid and ask price is called the ``spread'' and it
is the market maker's profit. The wider the spread, the more money the market maker makes.

\bd[Spread]
The \textbf{spread} is the difference between the ask price and the bid price:
\bse
\text{Spread} = \text{Ask Price} - \text{Bid Price}
\ese
\ed

The middle point between bid and ask prices is called the ``market price''.

\bd[Market Price]
The \textbf{market price} is the middle point between the bid and ask prices:
\bse
\text{Market Price} = \frac{\text{Bid Price} + \text{Ask Price}}{2}
\ese
\ed

The market price is the price at which a security is currently being traded in the secondary market. However, at the
market price, no one is able to neither buy nor sell a security. The bid and ask prices are the net prices at which a
market maker is willing to trade a security and that's the price that an investor will pay to buy a security or receive
to sell a security. \v

The difference between the market price and the bid and ask prices is called the ``mark-up'' and ``mark-down''.

\bd[Mark-Up]
A \textbf{mark-up} is the amount by which the ask price exceeds the market price:
\bse
\text{Mark-Up} = \text{Ask Price} - \text{Market Price}
\ese
\ed

\bd[Mark-Down]
A \textbf{mark-down} is the amount by which the bid price is less than the market price:
\bse
\text{Mark-Down} = \text{Market Price} - \text{Bid Price}
\ese
\ed

Saying the same thing in different words, market makers are eager to sell a security to an investor at market price
plus the mark-up, or buy a security from an investor at market price minus the mark-down. \v

One can easily show that the spread is equal to the sum of the mark-up and the mark-down:
\bse
\text{Spread} = \text{Mark-Up} + \text{Mark-Down}
\ese

Moving on, the second role of financial intermediaries in the secondary market is to act as brokers.

\bd[Broker]
A \textbf{broker} is a financial intermediary that brings buyers and sellers together but does not maintain an
inventory of securities.
\ed

A broker does not have capital at risk, and does not make money from the spread. Instead, brokers make money by
charging a commission for their services including compensation for time, effort and expenses of executing trade. The
commission is a fee that is charged for executing a trade. The commission can be a fixed amount, a percentage of the
trade, or a combination of both. \v

Secondary markets and their financial intermediaries provide several really important functions with the most important
one being liquidity. The fact that market makers make sure that securities can be bought and sold in the secondary
market provides investors with the ability to sell their securities and convert them into cash. This is important
because it means that investors can invest in securities without having to hold them until they mature. \v

Another important function of secondary markets is price discovery. The price of a security in the secondary market
provides information about the value of the security through the supply and demand for the security. This is important
because primary markets price things only once, at the time of the offering, and the price is set by the issuer,
without any input from the market and without having the possibility to change it, in a continuous changing world.

\subsection{Debt Securities / Fixed-Income Securities}

\bd[Debt Security / Fixed-Income Security]
A \textbf{debt security}, or \textbf{fixed-income security}, is a security that represents a creditor relationship with
a governmental body or a corporation.
\ed

In simple words, fixed-income securities are financial instruments that represent a loan made by an investor to the
issuer, making the investo a creditor of the issuer. The issuer of a fixed-income security is obligated to pay the
investor interest (hence the name ``fixed-income'') and to repay the principal amount at a later date, as it is
specified in the terms of repayment of the security. \v

\bd[Debt Markets / Fixed-Income Markets]
A \textbf{debt market}, or \textbf{fixed-income market}, is a market where people trade debt securities.
\ed

Fixed-income markets are usually double or triple the size of equity markets, and are considered to be the most
important type of securities market. \v

There are many different types of debt securities, but by far, the most important ones are the so-called ``bonds''.
\footnote{Bonds are so important that the terms ``debt security'', ``fixed-income security'', and ``bond'' are often
used interchangeably, despite the fact that technically they are not the same thing.}

\bd[Bond]
A \textbf{bond} is a fixed-income security that represents a loan made by an investor to a goverment, a govermental
agency or a corporate.
\ed

Bonds are the most important type of fixed-income securities. They are issued by governments, govermental agencies
and corporations to raise capital. Government and govermental agencies bonds (also called ``sovereign debts'') are
usually the largest and most liquid part of the fixed-income market, and are considered to be the safest type of
fixed-income securities. As a result of this, they have modest expected returns. Corporate bonds are much more complex
and difficult to categorize and analyze, however, they are usually riskier than government and govermental agencies
bonds, and as a result of this, they have higher expected returns. \v

When an investor buys a bond, they are lending money to the issuer in exchange for interest payments and the return of
the bond's principal (or face value) when it matures.

\bd[Principal / Face Value / Par Value]
The \textbf{principal}, \textbf{face value}, or \textbf{par value} of a bond is the amount that the issuer agrees to
repay the bondholder at the maturity date.
\ed

\bd[Maturity Date]
The \textbf{maturity date} is the date on which the issuer of a bond agrees to repay the bondholder the face value of
the bond.
\ed

The principal of the bond is usually paied out at its entirety at the end of the maturity date, however, there are
bonds that pay out the principal in installments, called ``amortizing bonds''.

\bd[Amortizing Bond]
An \textbf{amortizing bond} is a bond that pays out the principal in installments.
\ed

The interest payments on the principal that the issuer agrees to pay to the bondholder are called ``coupons''.

\bd[Coupon]
The \textbf{coupon} is the interest payment on the principal that the issuer agrees to pay to the bondholder at regular
intervals.
\ed

The coupons are usually made at regular intervals, e.g.\ annually, semi-annually, quarterly, or monthly, depending on
the fixed-income market. In this case the coupon is called ``periodic'', contrary to the case where the coupon is paied
at the end of the maturity date, called ``zero-coupon''.

\bd[Periodic Coupon]
A \textbf{periodic coupon} is a coupon that the issuer agrees to pay to the bondholder at regular intervals.
\ed

\bd[Zero-Coupon]
A \textbf{zero-coupon} is a coupon that is paied at its entirety at the end of the maturity date.
\ed

The interest rate that the issuer agrees to pay to the bondholder is called the ``coupon rate''.

\bd[Coupon Rate]
The \textbf{coupon rate} is the interest rate that the issuer agrees to pay to the bondholder.
\ed

The coupon rate is usually fixed, but it can also be floating, i.e.\ it can change over time.

\bd[Floating Coupon Rate]
A \textbf{floating coupon} is a coupon rate that can change over time.
\ed

The coupon rate is usually expressed as a percentage of the bond's face value, hence the coupons that the bondholder
receives are calculated by multiplying the coupon rate by the face value of the bond. \v

As it makes sense, bonds carry a certain amount of risks that makes the future return of the bondholder uncertain:
\bit
\item \textbf{Default / Credit Risk}: The risk that the issuer will not be able to make the interest payments or repay
the principal at the maturity date.
\item \textbf{Interest Rate Risk}: The risk that the value of the bond will decrease due to changes in interest rates.
This risk can be spllited into two subcategories: ``price risk'' and ``reinvestment risk''.
\bit
\item \textbf{Market / Price Risk}: The risk that the value of the bond will decrease due to an increase in interest
rates.
\item \textbf{Reinvestment Risk}: The risk that the bondholder will not be able to reinvest the interest payments at
the same interest rate.
\eit
\item \textbf{Inflation Risk}: The risk that the value of the bond will decrease due to an increase in inflation.
\item \textbf{Liquidity Risk}: The risk that the bondholder will not be able to sell the bond at a fair price.
\eit

Of course, since bonds carry risks they also carry sources of returns:
\bit
\item \textbf{Cpupons}: The interest payments that the bondholder receives from the issuer.
\item \textbf{Principal Capital Gains}: The increase in the face value of the bond due to a decrease in interest rates.
\item \textbf{Compound Interest}: The interest that the bondholder receives from reinvesting the interest payments.
\eit

\subsection{Equity Securities}

\bd[Equity Security]
An \textbf{equity security} is a security that represents ownership in a corporation.
\ed

Equity securities are also called ``stocks'' or ``shares''. When you buy a stock you are buying a piece of the
company. The company is divided into a number of shares and each share represents a piece of the company. When you
buy a share you become a shareholder of the company. The company is legally required to provide you with certain
rights and privileges as a shareholder. These rights and privileges include the right to vote in the company's annual
general meeting, the right to receive dividends, the right to receive a portion of the company's assets in case the
company is liquidated, and others. \v

\section{Derivatives}

\bd[Derivative]
A \textbf{derivative} is a contract for a future transacion between two parties, that derives its value from the
performance of an underlying entity.
\ed

\bd[Underlying]
The \textbf{underlying} is an asset, index, interest rate, or any variable from which a derivative is derived from.
\ed

Derivatives can be either ``exchange-traded'' or ``over-the-counter''.

\bd[Exchange-Traded Derivative]
An \textbf{exchange-traded derivative} is a standardized derivative that is traded on a regulated derivative exchange.
\ed

\bd[Derivative Exchange]
A \textbf{derivative exchange} is a central financial exchange where people can trade standardized exchange-traded
derivatives defined by the exchange.
\ed

Once two traders have agreed to a trade offered by an exchange, it is handled by the clearing house.

\bd[Clearing House]
A \textbf{clearing house} is an organization that facilitates tradings between two parties and manages the risks.
\ed

The advantage of a clearing house is that traders do not have to worry about the creditworthiness of the people they
are trading with. The clearing house takes care of credit risk by requiring each of the two traders to deposit funds
to ensure that they will live up to their obligations. \v

Traditionally exchange-traded derivatives traders used to physically meet on the floor of the exchange, shouting, and
using a complicated set of hand signals to indicate the trades they would like to carry out. Exchanges have largely
replaced this system by electronic trading which involves the use of algorithms to initiate trades, often without human
intervention, and has become an important feature of derivatives markets.

\bd[Over-The-Counter Derivative]
An \textbf{over-the-counter (OTC) derivative} is a derivative that is traded off-exchange.
\ed

Banks, other large financial institutions, fund managers, and corporations are the main participants in OTC derivatives
markets. \v

Once an OTC trade has been agreed, the two parties can either present it to a central counterparty (CCP), which is like
an exchange clearing house, or clear the trade bilaterally where the two parties have usually signed an agreement
covering all their transactions with each other.

\bd[Central Counterparty]
A \textbf{central counterparty (CCP)} stands between the two parties to the derivatives transaction so that one party
does not have to bear the risk that the other party will default.
\ed

While OTC markets were largely unregulated, following the financial crisis of 2008 we have seen the development of
many new regulations affecting the operation of OTC markets. The main objectives of the regulations are to improve
the transparency of OTC markets and reduce systemic risk. The over-the-counter market in some respects is being forced
to become more like the exchange-traded market. \v

Both OTC and exchange-traded markets for derivatives are huge. The number of derivatives transactions per year in OTC
markets is smaller than in exchange-traded markets, but the average size of the transactions is much greater, making the
volume of business in OTC market much larger than in the exchange-traded one. \v

Derivatives markets have been outstandingly successful. The main reason is that they have attracted many different
types of traders and have a great deal of liquidity. When a trader wants to take one side of a contract,
there is usually no problem in finding someone who is prepared to take the other side. \v

Three broad categories of traders can be identified: ``hedgers'', ``speculators'', and ``arbitrageurs''.

\bd[Hedger]
A \textbf{hedger} is a trader who uses derivatives to reduce the risk that they face from potential future movements in
a market variable.
\ed

\bd[Speculator]
A \textbf{speculator} is a trader who uses derivatives to bet on the future direction of a market variable.
\ed

\bd[Arbitrageur]
An \textbf{arbitrageur} is a trader who takes offsetting positions in two or more instruments to lock in a profit.
\ed

The fact that derivatives can be used for hedging, speculation, and arbitrage, makes them very versatile instruments.
Sometimes traders who have a mandate to hedge risks or follow an arbitrage strategy become (consciously or
unconsciously) speculators. The results can be disastrous, so it is very important for both financial and
nonfinancial corporations to set up controls to ensure that derivatives are being used for their intended purpose.
Risk limits should be set and the activities of traders should be monitored daily to ensure that these risk limits
are adhered to.

Unfortunately, even when traders follow the risk limits that have been specified, big mistakes can happen. Some of
the activities of traders in the derivatives market during the period leading up to the start of the financial crisis
in July 2007 proved to be much riskier than they were thought to be by the financial institutions they worked for.

\subsubsection{Forwards \& Futures}

\bd[Forward / Future]
A \textbf{forward} or \textbf{future} is a contract between two parties to buy or sell an underlying asset at a
specified price at a specified future date.
\ed

Forwards are traded in OTC markets, usually between two financial institutions or between a financial institution and
one of its clients. Futures are traded on exchanges, and the exchange acts as the counterparty to all trades.

\bd[Long Position]
The party that agrees to buy the asset is said to be in a \textbf{long position}.
\ed

\bd[Short Position]
The party that agrees to sell the asset is said to be in a \textbf{short position}.
\ed

\bd[Delivery Date / Maturity Date]
The specified future date $T$ in a forward or future is called the \textbf{delivery date} or \textbf{maturity date}.
\ed

\bd[Delivery Price / Forward Price / Future Price]
The specified price $K$ in a forward or future is called the \textbf{delivery price}, or \textbf{forward price}, or
\textbf{future price}.
\ed

Indepentendly of the delivery price, the underlying asset has its own price at each time step, called the ``spot
price''.

\bd[Spot Price]
The \textbf{spot price} $S_t$ is the price of an asset for immediate delivery at time t.
\ed

The payoff from a long position in a forward or future on one unit of an asset is:
\bse
S_T - K
\ese

where $K$ is the delivery price and $S_T$ is the spot price of the asset at maturity of the contract. This is because
the holder of the contract is obligated to buy an asset worth $S_T$ for $K$.

Similarly, the payoff from a short position in a forward of future contract on one unit of an asset is:
\bse
K - S_T
\ese

These payoffs can be positive or negative. Because it costs nothing to enter into a forward or future contract, the
payoff from the contract is also the trader's total gain or loss from the contract.

\vspace{5pt}
\fig{dv1}{0.65}

\subsubsection{Options}

\bd[Option]
An \textbf{option} is a contract that gives the holder the right, but not the obligation, to buy or sell an underlying
asset at a specified price at a specified future date.
\ed

It should be emphasized that an option gives the holder the right to do something. The holder does not have to
exercise this right. This is what distinguishes options from forwards and futures, where the holder is obligated to
buy or sell the underlying asset.

\bd[Strike Price / Exercise Price]
The specified price $K$ in an option is called the \textbf{strike price} or \textbf{exercise price}.
\ed

\bd[Expiration Date / Maturity]
The specified future date $T$ in an option is called the \textbf{expiration date} or \textbf{maturity}.
\ed

Options are traded both on exchanges and in the over-the-counter market. Based on if the option can be exercised only
on the expiration date or at any time up to the expiration date, we distinguish between ``European'' and ``American''
options.

\bd[European Option]
A \textbf{European option} is an option that can be exercised only on the expiration date.
\ed

\bd[American Option]
An \textbf{American option} is an option that can be exercised at any time up to the expiration date.
\ed

Note that the terms American and European do not refer to the location of the option or the exchange. Some options
trading on North American exchanges are European. Most of the options that are traded on exchanges are American.
European options are generally easier to analyze than American options, and some of the properties of an American
option are frequently deduced from those of its European counterpart. \v

There are two types of options: ``calls'' and ``puts''.

\bd[Call Option]
A \textbf{call option} is an option that gives the holder the right to buy an underlying asset at a specified price
at a specified future date.
\ed

\bd[Put Option]
A \textbf{put option} is an option that gives the holder the right to sell an underlying asset at a specified price
at a specified future date.
\ed

At this stage we note that there are four types of participants in options markets:
\bit
\item Buyers of calls.
\item Sellers of calls.
\item Buyers of puts.
\item Sellers of puts.
\eit

As with forwards and futures, buyers are referred to as having long positions and sellers are referred to as having
short positions. Selling an option is also known as ``writing'' the option. \v

Whereas it costs nothing to enter into a forward or futures contract, there is a cost to acquiring an option.]